% 11_board_io
\section{板级 I/O 与外设驱动}

板级 I/O 与外设驱动是 StarryOS 直接驱动 RK3588 板载外设的一层。SoC 把 UART、USB 主机控制器、PWM、网络控制器与引脚复用等外设都映射为内存中的寄存器块，这一层通过 MMIO 对这些寄存器块编程，并以设备节点与 sysfs 属性把它们交给用户态，从而为机器人的感知取流与执行器控制提供硬件通路，也为上板调试提供控制台、网络与重启原语。本报告中每一项在真机上测得的数据，都以这一层能在 RK3588 上稳定驱动外设为前提。两条板级现实决定了这一层的工作重心。其一，StarryOS 的用户 shell 一旦占用串口 tty，内核的 \texttt{warn!} 与 \texttt{info!} 便不再抵达控制台，\texttt{/dev/kmsg} 又仅可写入，第二条登录通路（SSH）因此成为硬性前提。其二，经 1.5 Mbaud 串口以 ymodem 刷写约 14 MB 的 FIT 映像长期不可靠，反复出现 \texttt{Error: NAK} 损坏，一夜运行中 8 次尝试全部失败，而任何一次物理下电都会重新枚举 USB 串口适配器，并可能在 8 核启动时导致电源欠压。两者共同要求一套以网络为先、暖重启为主的上板流程，这套流程又依赖可用的 \texttt{reboot} 系统调用、UART 与 USB 串口，以及板载网络。本节先说明这一层的构成与既有骨架，再逐组说明各外设驱动的机制以及既有与新增的分界，最后给出使能效果与刷写通路的实测。

\subsection{子系统构成与既有骨架}

StarryOS 与 ArceOS 此前已提供这一层的通用骨架。tty/pty 子系统提供字符设备到串口硬件之间的行规程（line discipline），负责规范模式下的行缓冲、回显与信号字符，并以等待队列在缓冲区为空时挂起读者、在字节到达时将其唤醒。USB 主机栈与 UVC 类驱动框架负责设备枚举、端点管理与传输请求的提交回收。sysfs 提供以属性文件暴露内核对象的通用模型。ax-net 提供网络协议栈与套接字接口。pinctrl 负责按引脚号把功能选择写入 IOMUX 寄存器。本队在此骨架上\ours{补齐 RK3588 的具体外设驱动}，并修复其中若干时序、锁与引脚路由缺陷，使骨架在真机上可用。整层按功能分为四组：控制台与执行器、双系统切换的重启原语、相机取流前端、从网卡上电到 SSH 的使能链。以下逐组说明，并标明各组中既有骨架与新增实现的分界。相关工作归本队板级侧，除引脚路由外均已在 \texttt{rcore-os/tgoskits} 合入。

\subsection{控制台与串口：UART 与 USB 串口 tty}

串口控制台是其余一切上板操作的前提。RK3588 的 UART 是一组 \texttt{DW\_apb\_uart} 兼容的 MMIO 寄存器块，含发送保持寄存器、接收缓冲寄存器、线路状态寄存器、FIFO 控制与波特率分频寄存器。RK3588 UART 初始化（\#1921）完成串口资源的建立：映射寄存器窗口、按源时钟设置分频以得到目标波特率、配置 FIFO 与线路控制、开放收发。这一步构成串口控制台的地基，此前缺少对 RK3588 UART 资源的初始化，控制台无从建立。

调试 UART 之外，第二条控制台通路来自 USB 串口。USB 串口 tty（\#1378）在 USB 主机栈枚举出串口适配器后将其绑定，以批量 IN/OUT 端点承载串行字节流，并把它作为一个 tty 暴露为 \texttt{/dev/ttyUSB} 设备。这条通路承担两件事：其一是提供第二个控制台，其二是机器人机械臂的串口控制器经此接入，机械臂在 \texttt{/dev/ttyUSB} 上呈现为一个 tty，用户态经该 tty 下发臂的动作指令。

交互式 shell 的可用性取决于 tty 输入路径的正确。TTY 输入刷新与读者唤醒修复（\#1922）补齐了 \texttt{TCFLSH} 对输入队列的刷新，并\ours{消除了一处丢失唤醒}：阻塞在 tty 上的读者在缓冲区为空时睡在等待队列上，输入路径每收到一个字节都应将其唤醒，此前某些情形下字节已入队而唤醒未发出，读者在尚有待读数据时继续睡眠，shell 因此挂起。修复后唤醒与数据到达严格对应，交互式 shell 自此可用。以上三项均已合并。

\subsection{执行器控制：PWM 的 sysfs 接口}

执行机构的控制面来自 PWM。PWM 以可编程的周期与占空比生成周期性方波，输出的平均电平由占空比决定，据此可设定舵机转角、电机转速与风扇转速。RK3588 的 PWM 控制器是一组 MMIO 寄存器块，含周期寄存器、占空比寄存器与控制寄存器，由源时钟驱动，编程时写入以时钟周期计的周期计数与占空比计数。\ours{PWM 的 sysfs 接口}（\#1468）按 Linux 的惯例在 \texttt{/sys/class/pwm} 下以 \texttt{period}、\texttt{duty\_cycle}、\texttt{enable}、\texttt{polarity} 等属性文件暴露该控制器，用户态写入这些文件即可驱动执行器。这是机器人底盘电机与舵机的支撑点，已合并。

\subsection{双系统切换：reboot 系统调用}

暖重启回路依赖一个可用的重启原语。\texttt{reboot(2)}（\#1358）在系统调用层\ours{实现 Linux 语义的 \texttt{reboot}}：调用方传入两个魔数与一个命令码（重启、关机、下电等），内核据此触发系统复位。在 RK3588 上，复位经固件的 PSCI \texttt{SYSTEM\_RESET} 调用落到硬件。有了这一原语，StarryOS shell 中一条 \texttt{reboot -f} 即可复位开发板并进入下一启动阶段，物理按键不再必要。这是暖重启回路成立的前提：StarryOS 经 \texttt{reboot} 进入 Linux，Linux 经 SSH 触发 \texttt{reboot} 进入 U-Boot，U-Boot 以 \texttt{fatload} 载入预置的 FIT 映像，两系统之间的切换全程不触碰电源按键。暖重启之所以关键，在于前述物理下电对 USB 串口重新枚举与 8 核上电欠压的影响。该系统调用已合并。

\subsection{视觉取流前端：UVC 相机异步传输}

USB 相机是整条视觉流水线的源头。UVC 是 USB 视频类标准，相机的帧经 USB 端点以流式方式传入：主机向端点提交一批传输请求，硬件在帧数据到达时逐一填充这些请求，完成回调将数据上交并立即重新提交请求以保持管道持续满载。这一提交、完成、重提交的循环即异步传输生命周期。UVC 异步传输生命周期修复（\#1924）\ours{纠正了这一循环中的时序错误}：生命周期处理不当会使已完成的请求未被重新提交，或在停止取流时状态错乱，取流因此停止或数据错乱。相机既是流水线的源头，这一修复对整条网球流水线直接承重。其后将 UVC 测试迁移至 \texttt{ax-sync}（\#1961），使相机测试路径改用正确的同步原语，消除测试中原有的同步语义错误。两项均已合并。

\subsection{从网卡上电到板上 SSH}

板上 SSH 取代了此前对单一串口通道的依赖，是上板迭代调试得以进行的关键。SSH 的成立取决于一条使能链，链上任一环缺失均无法登录。

链的底端是引脚路由。pinctrl 按引脚号将功能选择写入 IOMUX 寄存器，每个引脚的复用由 GRF/IOC 寄存器窗口中的若干比特决定，驱动需据引脚号与所在 bank 算出目标寄存器与位偏移。RK3588 的 pinctrl 原先误算了 IOMUX 的寄存器窗口偏移，Starry 写入了错误的窗口，GPIO、regulator 与外设的引脚复用因此错位。修正偏移计算后引脚复用方才正确，这一步是使能链的前置条件，现已实现并在板上验证，正作为独立 PR 整理中。

其上是网卡的上电、枚举与地址分配，由 rtnetlink 一组改动（\#1497）整合到位。网卡的电源轨由一个固定电压 regulator 门控，其使能是一个 GPIO，设备树声明该 GPIO 为 active-low，驱动须将其驱动至断言电平方能给电源轨供电，极性写反则网卡不上电；该组改动\ours{修正了固定 regulator 的 active-low 极性}，使电源轨上电。其次，板上有多颗 RTL8125 2.5G 网卡，多网卡枚举将每一颗暴露为一个网络接口。最后，rtnetlink 是内核与用户态之间基于 netlink 的网络配置接口，实现其 IPv4 地址操作后，用户态得以经 \texttt{ip addr} 为接口分配地址。三者串起，电源轨上电、端口枚举、地址分配依次到位，配合 TTY 修复（\#1922），脆弱的单一串口通道被替换为真正的登录，从网卡上电到 SSH 登录的整条链路在板上闭合。

网络路径另有一处正确性修复。ax-net 数据报锁修复（\#1963）避免了在数据报自旋锁下对互斥锁执行睡眠：持自旋锁时睡眠在原子上下文中不合法，负载下会引发错误睡眠，修复令数据报路径不再在持锁时睡眠于互斥锁。除引脚路由外，本组改动均已合并。

\subsection{使能效果与刷写通路}

这一层交付的是板级与外设的 bring-up，其效果集中体现在两处。其一，板上 SSH 登录取代了此前对单一串口的长期依赖，配合 \texttt{reboot} 原语构成不触碰电源按键的暖重启回路。其二，映像刷写从不可靠的串口 ymodem 转为以网络为先的通路。表~\ref{tab:board-flash} 给出三条刷写通路的实测用时。

\begin{table}[htbp]\centering
\caption{三条映像刷写通路的实测用时}\label{tab:board-flash}
\begin{tabular}{@{}p{4.6cm}p{2.6cm}p{6.2cm}@{}}
\toprule
通路 & 用时 & 说明 \\
\midrule
网络 \texttt{scp} FIT 至 \texttt{/boot} & 约 1 s & md5 校验一致，替代失败的 ymodem \\
串口 ymodem（1.5 Mbaud，约 14 MB FIT） & 单次 90 s 以上，含重试约 5 分钟 & 间歇 \texttt{Error: NAK}，一夜 8 次尝试全部失败 \\
U-Boot \texttt{fatload} 预置 FIT 后 \texttt{bootm} & 约 1.4 s & 相较 ymodem 的约 5 分钟 \\
\bottomrule
\end{tabular}
\end{table}

网络 \texttt{scp} 至 \texttt{/boot} 再经 \texttt{fatload} 与 \texttt{bootm} 启动的通路已在板上验证可用：载入修复后的内核，md5 校验一致并完成自检。这一层不含对 Linux 的性能主张，其内容为板级与外设的 bring-up。相机取流（\#1924/\#1961）与 PWM 控制面共同构成机器人的感知与执行通路，感知前端的模型结构优化与首次推理时延各自单独成章。本层的定位是把 RK3588 开发板变为可 SSH、可暖重启的实用目标，其余章节相对 Linux 的持平与领先方由此得以在真机上被测得。
